\section{Experiments}
\label{sec:experiments}

We evaluate \system{} by answering four questions.  First, can heterogeneous
execution sustain the 30-FPS real-time target once communication and runtime
overheads are included?  Second, does DAQ preserve metric-depth quality on
challenging sensor data?  Third, how does the complete system compare with
recent depth-completion methods?  Finally, what resource, communication,
power, and energy costs are required to obtain these gains?

\subsection{Settings}
\label{sec:settings}

\paragraph{Baselines}
The controlled system comparison runs the same depth-completion network in
three deployments: \emph{GPU-only} on an RTX~4090, \emph{Thor-only} on the
target edge SoC, and \system{}.  The GPU-only and Thor-only baselines execute
the complete model sequentially on one device; \system{} includes feature
transfer and output reconstruction in its end-to-end time.  For model-level
context, we additionally compare against CDM, PriorDA, PromptDA, and OMNI-DC,
using their released checkpoints and official evaluation procedures whenever
available.  Because their published FPS values come from different GPU
platforms, we label the execution platform of every method and do not treat
cross-platform FPS as a controlled speedup claim.

\paragraph{Benchmarks}
We evaluate all accuracy-oriented experiments on HAMMER and ClearPose.  These
datasets exercise different object geometries and difficult sensing conditions,
so results are reported for each dataset rather than averaged into one score.
Every method uses the same test split, validity mask, crop, input resolution,
and depth range within a dataset.  The system and ablation experiments use
batch size one, which matches continuous RGB-D capture on the target robotic
platform.

\paragraph{Experiment platform}
The prototype comprises an NVIDIA Jetson AGX Thor and an Intel Stratix FPGA
connected through the system interface described in
Section~\ref{sec:scheduling}.  The FPGA clock, external-memory configuration,
link rate, Thor power mode, and software versions are held fixed across all
\system{} variants.  The RTX~4090 baseline uses the same model weights, input
resolution, preprocessing, and output validation.

\paragraph{Metrics and measurement protocol}
We measure both single-frame latency and steady-state throughput.  Latency is
decomposed into preprocessing, encoder execution, boundary-tensor transfer,
decoding, return transfer, and output reconstruction; FPS is measured only
after the asynchronous pipeline is filled.  Timing includes synchronization
needed for a valid completed-depth output, rather than reporting kernel time
alone.  Depth quality is evaluated with AbsRel, RMSE, MAE, and
$\delta_{1.05}$.  Lower is better for the three error metrics and higher is
better for $\delta_{1.05}$.  We report system power, energy per completed
frame, and FPS/W under a consistent measurement boundary.

\subsection{Main Experiments}
\label{sec:main-exp}

% Main results (table*): its source stays on page 4 so it lands on top of
% page 5.
\begin{table*}[t]
  \caption{Main results. Values marked ``--'' await measurement.}
  \label{tab:main-exp}
  \centering
  \small
  \setlength{\tabcolsep}{4pt}
  \renewcommand{\arraystretch}{1.10}
  \begin{tabular}{llc|ccc|ccc}
    \toprule
    & & & \multicolumn{3}{c|}{HAMMER} & \multicolumn{3}{c}{ClearPose} \\
    Method & Platform & FPS $\uparrow$
      & AbsRel $\downarrow$ & RMSE $\downarrow$ & $\delta_{1.05}$ $\uparrow$
      & AbsRel $\downarrow$ & RMSE $\downarrow$ & $\delta_{1.05}$ $\uparrow$ \\
    \midrule
    CDM & -- & -- & -- & -- & -- & -- & -- & -- \\
    PriorDA & -- & -- & -- & -- & -- & -- & -- & -- \\
    PromptDA & -- & -- & -- & -- & -- & -- & -- & -- \\
    OMNI-DC & -- & -- & -- & -- & -- & -- & -- & -- \\
    \system{} & -- & -- & -- & -- & -- & -- & -- & -- \\
    \bottomrule
  \end{tabular}
\end{table*}

\paragraph{Throughput}
Table~\ref{tab:main-exp} compares the inference rate of CDM, PriorDA,
PromptDA, OMNI-DC, and \system{}, together with the execution platform of
each method.  This comparison answers a deployment question---whether current
depth-completion methods can service a real-time sensor stream---while the
controlled GPU-only/Thor-only comparison supports architectural speedup
claims.  We report \system{}'s measured end-to-end FPS and identify whether its
advantage comes from the coarse model partition, cross-frame overlap, or
low-precision execution, using the ablation evidence rather than the
cross-platform numbers alone.

\paragraph{Depth accuracy}
Table~\ref{tab:main-exp} also compares all methods under common HAMMER and
ClearPose protocols.  The discussion reports the absolute metric for every
method, \system{}'s gap to the best-performing method, and its rank on each
dataset.  The relevant question is whether \system{} reaches real-time edge
throughput while keeping the accuracy loss small and consistent across both
datasets.

\subsection{Ablation Studies}
\label{sec:ablation}

\paragraph{Optical flow versus frequency}
\draftnote{Ablation on optical flow at different rates, shown in
Fig.~\ref{fig:of-vs-hz}.}

% OF vs. Hz figure: bottom of the left column on page 5.
\begin{figure}[b]
  \centering
  \resultslot{1.30in}{Optical flow vs.\ Hz}
  \caption{Ablation on optical flow versus frequency. \pending}
  \label{fig:of-vs-hz}
\end{figure}

\paragraph{Scheduling}
\draftnote{Ablation on scheduling, reported in Table~\ref{tab:abl-scheduling}.}

\paragraph{Heterogeneous execution}
Thor-only versus \system{} is the primary controlled comparison for the
heterogeneous partition: both process the same frames with the same network,
while only the latter offloads the decoder and overlaps encoder,
communication, and decoder work across frames.  We therefore report not only
their FPS difference, but also the stage-latency breakdown and the fraction of
transfer hidden by the asynchronous schedule.  GPU-only is retained as a
high-end performance reference, not as an edge-device energy baseline.

\paragraph{Metric-depth accuracy}
The decisive accuracy comparison is \system{} without versus with DAQ because
their model partition, target bit widths, and runtime are identical.  Any
quality difference therefore measures the value of coordinating the quantized
encoder MLP, boundary representation, and decoder using metric-depth
supervision.  We report each dataset independently and discuss whether
improvements come from reducing large depth errors, increasing the fraction of
accurate pixels, or both.

% Temporary: the draft body does not fill the left column of page 5 yet, so
% end it here; Tables III and IV are then read in the right column. Remove
% once the ablation text fills the left column.
\newpage

% Scheduling ablation: right column on page 5, below Table II.
\begin{table}[t]
  \caption{Ablation on scheduling. \pending}
  \label{tab:abl-scheduling}
  \centering
  \tabslot{0.80in}{Ablation on scheduling}
\end{table}

% Fallback table: bottom of the right column on page 5.
\begin{table}[b]
  \caption{Fallback optical flow on CUDA. \pending}
  \label{tab:fallback}
  \centering
  \tabslot{0.80in}{Fallback optical flow on CUDA}
\end{table}

\paragraph{Fallback optical flow on CUDA}
\draftnote{Optical flow falls back to CUDA; results in Table~\ref{tab:fallback}.}

\paragraph{Power and energy efficiency}
Power is measured for Thor-only and the complete heterogeneous system using
the same input stream and steady-state window used for FPS.  We derive energy
per frame as system power divided by completed frames per second and
additionally report FPS/W.  The analysis separates the extra power introduced
by the additional hardware from the energy saved by shortening the end-to-end
execution time.

% Visualization (figure*): its source stays on page 5 so it lands on top of
% page 6.
\begin{figure*}[t]
  \centering
  \figslot{figures/visualization.pdf}{2.40in}{%
    Qualitative visualization with zoom-in regions}
  \caption{Visualization. \pending}
  \label{fig:visualization}
\end{figure*}

% Temporary: the draft body does not fill page 5 yet, so start page 6 here to
% keep Fig. 6 above Detailed Visualization. Remove once page 5 is full.
\newpage

\subsection{Detailed Visualization}
\label{sec:visualization}

\draftnote{Detailed visualization in Fig.~\ref{fig:visualization}, followed by
a zoom-in explanation.}
